\SourcePage{295}{298}
\section{Second quantization. The case of Bose statistics}

For the theory of systems formed by many identical particles, a special
method known as \emph{second quantization} is used extensively. This method is
particularly indispensable in relativistic theory, where

\SourcePage{296}{299}
one has to treat systems in which even the number of particles is
variable\footnote{Dirac developed the method of second quantization for
photons in the theory of radiation (1927); Wigner and Jordan subsequently
extended it to fermions (E.~Wigner, P.~Jordan, 1928).}.

Let $\psi_1(\xi)$, $\psi_2(\xi)$,~\dots be some complete orthogonal normalized
set of wave functions for the stationary states of one
particle\footnote{As in \S~61, $\xi$ denotes all the coordinates together
with the particle's spin projection $\sigma$; integration over $d\xi$ will
mean integration over the coordinates and summation over $\sigma$.}. These
may be the states of a particle in an arbitrarily selected external field,
although plane waves are generally chosen: the wave functions of a free
particle having specified momentum (and spin projection). To make the state
spectrum discrete, the motion is then considered in a large but bounded
region of space. For motion in a finite volume, the eigenvalues of the
momentum components form a discrete sequence; the spacings of successive
values are inversely proportional to the linear dimensions of the region and
tend to zero as those dimensions increase.

In a system of free particles, the momentum of every particle is conserved
separately. Consequently, the \emph{occupation numbers} are conserved as
well: the numbers $N_1,N_2,\ldots$ that state how many particles are in each
of the states $\psi_1,\psi_2,\ldots$. In a system of interacting particles,
the separate particle momenta are no longer conserved, and neither are the
occupation numbers. For such a system one can speak only of the probability
distribution among different values of the occupation numbers. Our aim is to
construct a mathematical formalism in which these numbers themselves, rather
than particle coordinates and spin projections, are the independent
variables.

Dirac notation (see the end of \S~11) is convenient in this formalism, with
$N_1,N_2,\ldots$ chosen as the quantum numbers specifying a state. The states
corresponding to the wave functions (61.3) and (61.5) will be denoted by
$|N_1,N_2,\ldots\rangle$. Coordinate and spin variables then no longer appear
explicitly.

With this choice of independent variables, the operators of the various
physical quantities (including the system's Hamiltonian) must likewise be
specified by how they act on functions of the occupation numbers. Such a
specification can be obtained by starting from the usual matrix representation
of operators. One must then consider

\SourcePage{297}{300}
the operator matrix elements between wave functions for stationary states of
the noninteracting-particle system. Since these states can be characterized by
assigning definite values to the occupation numbers, the resulting matrix
elements reveal how the operators act on those variables.


We first consider systems of particles obeying Bose statistics.

Let $\hat f_a^{(1)}$ be the operator of some quantity referring to one
particle (particle $a$), so that it acts only on functions of the variables
$\xi_a$. Introduce the operator symmetric in all particles
\begin{equation}
 \hat F^{(1)}=\sum_a\hat f_a^{(1)}
 \tag{64.1}
\end{equation}
(the sum runs over every particle), and determine its matrix elements with
respect to the wave functions (61.3). It is readily seen first of all that
these matrix elements can be nonzero only for transitions that leave all the
numbers $N_1,N_2,\ldots$ unchanged (the diagonal elements), or for transitions
in which one of these numbers rises by one while another falls by one. Indeed,
each operator $\hat f_a^{(1)}$ acts on only one factor in the product
$\psi_{p_1}(\xi_1)\psi_{p_2}(\xi_2)\cdots\psi_{p_N}(\xi_N)$, and hence its
matrix element can differ from zero only when a single particle changes its
state. This means that the number of particles in one state decreases by one,
whereas the number in another state increases by one. The actual evaluation
of these matrix elements is very simple; it is easier to carry it out than to
follow a written account of it. We therefore give only the result. The
off-diagonal elements are
\begin{equation}
 \langle N_i,N_k-1|\hat F^{(1)}|N_i-1,N_k\rangle
 =f_{ik}^{(1)}\sqrt{N_iN_k}.
 \tag{64.2}
\end{equation}
Only the indices in which the matrix element is not diagonal have been shown;
the others are suppressed for brevity. Here $f_{ik}^{(1)}$ is the matrix
element
\begin{equation}
 f_{ik}^{(1)}=\int\psi_i^*(\xi)\hat f^{(1)}\psi_k(\xi)\,d\xi;
 \tag{64.3}
\end{equation}
the operators $\hat f_a^{(1)}$ differ only in the designation of the variables
on which they act, so the integrals (64.3) do not depend on $a$, and this index
has been omitted. The diagonal matrix elements of $\hat F^{(1)}$ are the mean
values

\SourcePage{298}{301}
of $F^{(1)}$ in the states $\Psi_{N_1N_2\ldots}$. Their evaluation gives
\begin{equation}
 \overline{F^{(1)}}=\sum_i f_{ii}^{(1)}N_i.
 \tag{64.4}
\end{equation}

We now introduce the operators fundamental to the second-quantization method,
$\hat a_i$. They act not on functions of coordinates but on functions of the
occupation numbers. By definition, when $\hat a_i$ acts on the state
$|N_1,N_2,\ldots\rangle$, it reduces the variable $N_i$ by one and at the same
time multiplies the function by $\sqrt{N_i}$\footnote{We have used the natural
notation $\hat a|n\rangle$ for the result of applying the operator $\hat a$ to
the wave function of the state $|n\rangle$.}:
\begin{equation}
 \hat a_i|N_1,N_2,\ldots,N_i,\ldots\rangle
 =\sqrt{N_i}|N_1,N_2,\ldots,N_i-1,\ldots\rangle.
 \tag{64.5}
\end{equation}
One may say that $\hat a_i$ lowers by one the number of particles in the
$i$th state; it is accordingly called the particle \emph{annihilation
operator}. It can be represented by a matrix having just one nonzero element,
namely
\begin{equation}
 \langle N_i-1|\hat a_i|N_i\rangle=\sqrt{N_i}.
 \tag{64.6}
\end{equation}

The operator $\hat a_i^+$ adjoint to $\hat a_i$ is represented, by definition
(see (11.9)), by the matrix with the single element
\begin{equation}
 \langle N_i|\hat a_i^+|N_i-1\rangle
 =\langle N_i-1|\hat a_i|N_i\rangle^*=\sqrt{N_i}.
 \tag{64.7}
\end{equation}
This means that its action on the function $|N_1,N_2,\ldots\rangle$ increases
$N_i$ by one:
\begin{equation}
 \hat a_i^+|N_1,N_2,\ldots,N_i,\ldots\rangle
 =\sqrt{N_i+1}|N_1,N_2,\ldots,N_i+1,\ldots\rangle.
 \tag{64.8}
\end{equation}
In other words, $\hat a_i^+$ increases by one the number of particles in the
$i$th state; it is called the particle \emph{creation operator}.

When the product $\hat a_i^+\hat a_i$ acts on a wave function, it can only
multiply that function by a constant while leaving all the variables
$N_1,N_2,\ldots$ unchanged: $\hat a_i$ lowers $N_i$ by one, after which
$\hat a_i^+$ restores its initial value. Direct multiplication of the matrices
(64.6) and (64.7) does indeed show that $\hat a_i^+\hat a_i$ is represented
by a diagonal matrix whose diagonal entries equal $N_i$:
\begin{equation}
 \hat a_i^+\hat a_i=N_i.
 \tag{64.9}
\end{equation}
In the same way one finds
\begin{equation}
 \hat a_i\hat a_i^+=N_i+1.
 \tag{64.10}
\end{equation}

\SourcePage{299}{302}
Subtracting these expressions gives the commutation rule for $\hat a_i$ and
$\hat a_i^+$:
\begin{equation}
 \hat a_i\hat a_i^+-\hat a_i^+\hat a_i=1.
 \tag{64.11}
\end{equation}
Operators with different indices $i$ and $k$, which act on different
variables ($N_i$ and $N_k$), commute:
\begin{equation}
 \hat a_i\hat a_k-\hat a_k\hat a_i=0,\qquad
 \hat a_i\hat a_k^+-\hat a_k^+\hat a_i=0,\qquad i\ne k.
 \tag{64.12}
\end{equation}

From the properties of $\hat a_i$ and $\hat a_i^+$ just described, it follows
at once that the operator
\begin{equation}
 \hat F^{(1)}=\sum_{i,k}f_{ik}^{(1)}\hat a_i^+\hat a_k
 \tag{64.13}
\end{equation}
is identical with (64.1). Indeed, all matrix elements computed with (64.6)
and (64.7) coincide with those in (64.2) and (64.4). This result is very
important. The quantities $f_{ik}^{(1)}$ in (64.13) are merely numbers. We
have thereby expressed an ordinary operator acting on functions of the
coordinates as an operator acting on functions of the new variables, the
occupation numbers $N_i$.

The result is readily generalized to operators of other kinds. Let
\begin{equation}
 \hat F^{(2)}=\sum_{a>b}\hat f_{ab}^{(2)},
 \tag{64.14}
\end{equation}
where $\hat f_{ab}^{(2)}$ is the operator of a physical quantity referring to
a pair of particles at once and therefore acts on functions of $\xi_a$ and
$\xi_b$. An analogous calculation shows that this operator can be expressed
through $\hat a_i$ and $\hat a_i^+$ as
\begin{equation}
 \hat F^{(2)}=\frac12\sum_{i,k,l,m}
 \langle ik|f^{(2)}|lm\rangle\hat a_i^+\hat a_k^+\hat a_m\hat a_l,
 \tag{64.15}
\end{equation}
\[
 \langle ik|f^{(2)}|lm\rangle=\iint
 \psi_i^*(\xi_1)\psi_k^*(\xi_2)\hat f^{(2)}
 \psi_l(\xi_1)\psi_m(\xi_2)\,d\xi_1d\xi_2.
\]
The generalization of these formulas to operators of any other kind that are
symmetric in all the particles ($\hat F^{(3)}=\sum\hat f_{abc}^{(3)}$, and so
forth) is evident.

These formulas also allow the Hamiltonian of the physical system under study,
consisting of $N$ interacting identical particles, to be expressed through
$\hat a_i$ and $\hat a_i^+$. The Hamiltonian of such a system is, of course,
symmetric in all the particles. In the nonrelativistic
approximation\footnote{In the absence of a magnetic field.}, it is independent
of their spins

\SourcePage{300}{303}
and may be written in the general form
\begin{equation}
 \hat H=\sum_a\hat H_a^{(1)}+\sum_{a>b}U^{(2)}(\mathbf r_a,\mathbf r_b)
 +\sum_{a>b>c}U^{(3)}(\mathbf r_a,\mathbf r_b,\mathbf r_c)+\cdots
 \tag{64.16}
\end{equation}
Here $\hat H_a^{(1)}$ is the part of the Hamiltonian depending on the
coordinates of only one particle, particle $a$:
\begin{equation}
 \hat H_a^{(1)}=-\frac{\hbar^2}{2m}\Delta_a+U^{(1)}(\mathbf r_a),
 \tag{64.17}
\end{equation}
where $U^{(1)}(\mathbf r_a)$ is the potential energy of one particle in the
external field. The remaining terms in (64.16) describe the energy of
interaction among the particles; terms depending respectively on the
coordinates of two, three, and so on particles have been separated.

Writing the Hamiltonian in this form permits direct application of (64.13),
(64.15), and the analogous formulas. Thus
\begin{equation}
 \hat H=\sum_{i,k}H_{ik}^{(1)}\hat a_i^+\hat a_k
 +\frac12\sum_{i,k,l,m}\langle ik|U^{(2)}|lm\rangle
 \hat a_i^+\hat a_k^+\hat a_m\hat a_l+\cdots
 \tag{64.18}
\end{equation}
This yields the required representation of the Hamiltonian as an operator
acting on functions of the occupation numbers.


When the particles in a system do not interact, expression (64.18) retains only
its first term:

\begin{equation}
 \hat H=\sum_{i,k}H_{ik}^{(1)}\hat a_i^+\hat a_k.
 \tag{64.19}
\end{equation}
If the functions $\psi_i$ are chosen to be eigenfunctions of the
one-particle Hamiltonian $\hat H^{(1)}$, then the matrix $H_{ik}^{(1)}$ is
diagonal, and its diagonal entries are the particle's energy eigenvalues
$\varepsilon_i$. Consequently,
\[
 \hat H=\sum_i\varepsilon_i\hat a_i^+\hat a_i;
\]
replacing $\hat a_i^+\hat a_i$ by its eigenvalues (64.9), we obtain the
following expression for the energy levels of the system:
\[
 E=\sum_i\varepsilon_iN_i.
\]
This is the trivial result that should indeed have emerged.

\SourcePage{301}{304}
The formalism developed above can be put into a more compact form by
introducing the so-called $\psi$-operators\footnote{Notice the analogy between
(64.20) and the expansion $\psi=\sum a_i\psi_i$ of an arbitrary wave function
in a complete set of functions. Here the expansion is, as it were, quantized
once again, which accounts for the name of the entire method: second
quantization.}
\begin{equation}
 \hat\psi(\xi)=\sum_i\psi_i(\xi)\hat a_i,\qquad
 \hat\psi^+(\xi)=\sum_i\psi_i^*(\xi)\hat a_i^+,
 \tag{64.20}
\end{equation}
where $\xi$ is regarded as a parameter. It is clear from the properties of
$\hat a_i$ and $\hat a_i^+$ stated above that $\hat\psi$ decreases the total
number of particles in the system by one, whereas $\hat\psi^+$ increases it
by one.

It is easy to see that $\hat\psi^+(\xi_0)$ creates a particle situated at the
point $\xi_0$. Indeed, the action of $\hat a_i^+$ creates a particle in the
state whose wave function is $\psi_i(\xi)$. It follows that applying
$\psi_i^*(\xi_0)\hat a_i^+$ creates a particle in the state with wave function
\[
 \sum_i\psi_i^*(\xi)\psi_i(\xi_0)=\delta(\xi-\xi_0).
\]
Formula (5.12) has been used here; the result corresponds to a particle with
specified coordinates and spin\footnote{$\delta(\xi-\xi_0)$ conventionally
denotes the product
\[
 \delta(x-x_0)\delta(y-y_0)\delta(z-z_0)\delta_{\sigma\sigma_0}.
\]}.

The commutation rules for the $\psi$-operators follow immediately from those
for $\hat a_i$ and $\hat a_i^+$:
\begin{equation}
 \hat\psi(\xi)\hat\psi(\xi')-\hat\psi(\xi')\hat\psi(\xi)=0,
 \tag{64.21}
\end{equation}
\begin{equation}
 \hat\psi(\xi)\hat\psi^+(\xi')-\hat\psi^+(\xi')\hat\psi(\xi)
 =\sum_i\psi_i(\xi)\psi_i^*(\xi')=\delta(\xi-\xi').
 \tag{64.22}
\end{equation}

In terms of the $\psi$-operators, the second-quantized operator
$\hat F^{(1)}$ takes the form
\begin{equation}
 \hat F^{(1)}=\int\hat\psi^+(\xi)\hat f^{(1)}\hat\psi(\xi)\,d\xi
 \tag{64.23}
\end{equation}
(the operator $\hat f^{(1)}$ is understood here to act in $\hat\psi(\xi)$ on
functions of the parameters $\xi$). Substitution of $\hat\psi$ and
$\hat\psi^+$ from (64.20), together with the definition (64.3), does indeed
return us

\SourcePage{302}{305}
to (64.13). Similarly, in place of (64.15) we have
\begin{equation}
 \hat F^{(2)}=\frac12\iint\hat\psi^+(\xi)\hat\psi^+(\xi')
 \hat f^{(2)}\hat\psi(\xi')\hat\psi(\xi)\,d\xi\,d\xi'.
 \tag{64.24}
\end{equation}

In particular, the Hamiltonian of the system expressed through the
$\psi$-operators is
\begin{equation}
 \begin{aligned}
 \hat H={}&\int\left\{-\frac{\hbar^2}{2m}\hat\psi^+(\xi)
 \Delta\hat\psi(\xi)+\hat\psi^+(\xi)U^{(1)}(\xi)\hat\psi(\xi)\right\}d\xi\\
 &+\frac12\iint\hat\psi^+(\xi)\hat\psi^+(\xi')U^{(2)}(\xi,\xi')
 \hat\psi(\xi')\hat\psi(\xi)\,d\xi\,d\xi'+\cdots
 \end{aligned}
 \tag{64.25}
\end{equation}

The operator $\hat\psi^+(\xi)\hat\psi(\xi)$ is formed from the $\psi$-operators
in analogy with the product $\psi^*\psi$, which gives the probability density
of a particle in a state whose wave function is $\psi$. It is called the
particle \emph{density operator}. The integral

\begin{equation}
 \hat N=\int\hat\psi^+\hat\psi\,d\xi
 \tag{64.26}
\end{equation}
serves in the second-quantization formalism as the operator for the total
number of particles in the system. Indeed, on substituting the
$\psi$-operators from (64.20) and using the normalization and mutual
orthogonality of the wave functions, we find
$\hat N=\sum_i\hat a_i^+\hat a_i$. Each term in this sum is the particle-number
operator for the $i$th state; by (64.9), its eigenvalues are the occupation
numbers $N_i$. Their sum is the total number of particles in the
system\footnote{For systems having a prescribed number of particles, these
statements (as well as the properties of the free-particle Hamiltonian
(64.19)) appear trivial. In relativistic theory, however, their generalization
leads to new results which are far from trivial (cf. Vol.~IV, \S~11).}.


Finally, if a system contains bosons of different kinds, separate operators
$\hat a$ and $\hat a^+$ must be introduced for every kind of particle in the
second-quantization method. Operators belonging to different kinds of
particles evidently commute with one another.
